\subsection{Startovní prodleva a synchronizace s hráčem}

Před samotným pohybem běží korutina \texttt{StartMovingAfterDelay()}, která čeká \texttt{startDelay} sekund. Během této doby se \texttt{\_canMove} nastaví na false a pohyb se neprovádí. To dává hráči náskok před tím, než začne nebezpečí v podobě pronásledování koulí.

Tato prodleva je důležitá pro herní design, jelikož bez ní by se koule začala hned pohybovat ve směru hráče, což by mohlo působit nespravedlivě. Startovní delay poskytuje férový start a umožňuje hráči dostat se do bezpečnější vzdálenosti.

Stejná prodleva se restartuje při respawnu, aby po smrti měl hráč opět čas na přípravu. To zajišťuje konzistentní herní podmínky a zabraňuje situacím, kdy by se koule začala pohybovat k respawnujícímu se hráči okamžitě.